\section*{Procesja}

\gresetinitiallines{1}
\comm{Idąc w procesji do Najświętszej Maryi Panny, śpiewa się następującą antyfonę:}\\
\gregorioscore[a]{GABC/Procesja/salve.gabc}\hspace{1mm}\\
\gresetinitiallines{0}
\comm{Po antyfonie następuje dialog i oracja:}\\
\gregorioscore[a]{GABC/Procesja/salve_dialog.gabc}\hspace{1mm}\\
\gregorioscore[a]{GABC/Procesja/salve_oracja.gabc}\hspace{1mm}\\

\gresetinitiallines{1}
\comm{Wracając do chóru, śpiewa się następującą antyfonę do św. Dominika:}\\
\gregorioscore[a]{GABC/Procesja/lumen.gabc}\hspace{1mm}\\
\gresetinitiallines{0}
\comm{Po antyfonie następuje dialog i oracja:}\\
\gregorioscore[a]{GABC/Procesja/lumen_dialog.gabc}\hspace{1mm}\\
\gregorioscore[a]{GABC/Procesja/lumen_oracja.gabc}\hspace{1mm}\\

\gresetinitiallines{1}
\comm{Wedle zwyczaju prowincji Polskiej w środę po Salve zamiast "O Lumen" śpiewana jest antyfona do św. Jacka:}\\
\gregorioscore[a]{GABC/Procesja/ave_florum.gabc}\hspace{1mm}\\
\gresetinitiallines{0}
\comm{Po antyfonie następuje dialog i oracja:}\\
\gregorioscore[a]{GABC/Procesja/ave_florum_dialog.gabc}\hspace{1mm}\\
\gregorioscore[a]{GABC/Procesja/ave_florum_oracja.gabc}\hspace{1mm}\\